\begin{abstract}
  We study \emph{distribution-aware reliability estimation}: given a
  deterministic program $\prog$, a discrete distribution $\dist$ over
  its inputs, and a Boolean property $\prop$, certify a confidence
  interval for $\Pr_{x \sim \dist}[\prop(\prog(x)) = 1]$ using as few
  program evaluations as possible.  Rare-event programs defeat plain
  sampling, while exact symbolic techniques fail on programs with
  non-linear arithmetic, unbounded loops, or non-uniform input
  distributions.

  We present \toolname{}, an adaptive partitioning algorithm that
  bridges the two regimes.  \toolname{} maintains a frontier of
  path-condition regions and at each step selects between
  \textsc{Sample} and SMT-driven \textsc{Refine} under a unified
  gain-per-cost criterion.  The certified half-width decomposes as
  $\Pr\!\bigl[|\muhat - \mutrue| \le \epsstat + \Wopen\bigr] \ge
  1 - \delta$, where $\epsstat$ aggregates per-region statistical
  uncertainty and $\Wopen$ is the total probability mass of regions
  not yet symbolically certified.  This decomposition is
  simultaneously the algorithm's correctness theorem and its
  scheduler: \textsc{Sample} reduces $\epsstat$, \textsc{Refine}
  reduces $\Wopen$, and the algorithm terminates as soon as their sum
  drops below the target $\varepsilon$.

  The per-region statistical bound instantiates the predictable-plug-in
  empirical-Bernstein confidence sequence of Waudby-Smith and
  Ramdas---closed-form, variance-adaptive, and anytime-valid under
  \toolname{}'s data-dependent stopping rule.  The closure rule is
  single-branch: a region is certified if and only if an SMT oracle
  proves path-determinism, with no sample-based fallback, yielding
  a one-lemma soundness contract.

  We implement \toolname{} as a Python prototype with a Z3 backend
  and evaluate it on twelve integer benchmarks drawn from established
  sources (Hacker's Delight, CLRS, the CERT~C Secure Coding standard,
  the Collatz literature) plus a parametric rare-event program.  Two
  ablations decompose the contribution of the bound and of the
  stratification: holding the bound constant, \toolname{}'s symbolic
  stratification reaches the same half-width with $9\times$ fewer
  samples; holding the adaptive-allocation strategy constant,
  symbolic guidance is $4\times$ tighter than variance-driven random
  stratification.  On a parametric rare-event program, \toolname{}
  certifies $\mutrue$ to half-width exactly zero in $\Theta(1)$
  samples in the sweet-spot regime $\mutrue \in [5\times 10^{-4},
  2\times 10^{-3}]$---a regime where plain Monte Carlo's
  $\Theta(1/(\mutrue\varepsilon^2))$ sample complexity is prohibitive.
\end{abstract}
